\subsection{Substitutions}

Let E be a unital associative algebra over A and \(\mathbf{x} = (x_{i})_{i \in I}\) a family of pairwise permutable elements of E. Let \(\mathbf{X} = (\mathbf{X}_{i})_{i \in I}\) be a family of indeterminates. By III, Prop. 7, p.~449, there exists a unique unital homomorphism f of A{[}X{]} into E such that \(f(\mathbf{X}_{i}) = x_{i}\) for every \(i \in I\) . The image of an element u of A{[}X{]} under f is written \(u(\mathbf{x})\) and is called the element of E obtained by substituting \(x_{i}\) for \(X_{i}\) in u, or also the value of u for \(X_{i} = x_{i}\) . In particular, \(u = u((\mathbf{X}_{i})_{i \in I})\) . If \(I = \{1, \ldots, n\}\) we write \(u(x_{1}, \ldots, x_{n})\) in place of \(u((x_{i})_{i \in I})\) . More generally, if M is an A-module and if v is an element of

\[
\mathbf {M} [ (\mathbf {X} _ {i}) _ {i \in I} ] = \mathbf {M} \otimes_ {\mathbf {A}} \mathbf {A} [ (\mathbf {X} _ {i}) _ {i \in I} ],
\]

we denote the image of \(v\) in \(\mathbf{M} \otimes_{\mathbf{A}} \mathbf{E} = \mathbf{M}_{(\mathbf{E})}\) under the mapping \(1_{\mathbf{M}} \otimes f\) by \(v(\mathbf{x})\) .

If the homomorphism \(u \mapsto u(\mathbf{x})\) of \(\mathbf{A}[\mathbf{X}]\) into \(\mathbf{E}\) is injective, we say that the family \(\mathbf{x}\) is algebraically free over \(\mathbf{A}\) , or that the \(x_i\) are algebraically independent over \(\mathbf{A}\) . This also means that the monomials \(\mathbf{x}^\nu\) ( \(\nu \in \mathbf{N}^{(I)}\) ) are linearly independent over \(\mathbf{A}\) .

If \(\lambda\) is a unital homomorphism of \(E\) into a unital associative A-algebra \(E'\) , we have

\[
\lambda \left(u \left(\left(x _ {i}\right) _ {i \in I}\right)\right) = u \left(\left(\lambda \left(x _ {i}\right) _ {i \in I}\right)\right),\tag{1}
\]

because \(\lambda \circ f\) is a homomorphism of \(\mathbf{A}[\mathbf{X}]\) into \(\mathbf{E}'\) which maps \(\mathbf{X}_i\) to \(\lambda(x_i)\) .

Let \(u \in A[X]\) . If E is commutative, the mapping \(x \mapsto u(x)\) of \(E^{I}\) into E is called the polynomial function defined by \(u\) (and the algebra E); we shall sometimes denote it by \(\tilde{u}\) (or even just \(u\) ).

Let \(\mathbf{Y} = (\mathbf{Y}_j)_{j\in J}\) be another family of indeterminates, and let us take for E the polynomial algebra \(\mathbf{A}[\mathbf{Y}]\) . Given \(u\in \mathbf{A}[\mathbf{X}]\) , take \(g_{i}\in \mathbf{A}[\mathbf{Y}]\) for \(i\in I\) and put \(\mathbf{g} = (g_i)_{i\in I}\) ; let \(u(\mathbf{g})\in \mathbf{A}[\mathbf{Y}]\) be the polynomial obtained by substituting the polynomials \(g_{i}\) for \(\mathbf{X}_i\) in \(u\) . Let \(\mathbf{y} = (y_j)_{j\in J}\) be a family of pairwise permutable elements of a unital associative A-algebra \(\mathbf{E}'\) ; applying (1) and taking for \(\lambda\) the homomorphism \(g\mapsto g(\mathbf{y})\) of E into \(\mathbf{E}'\) , we obtain

\[
(u (\mathbf {g})) (\mathbf {y}) = u \left(\left(g _ {i} (\mathbf {y})\right)\right).\tag{2}
\]

If \(\mathbf{f} = (f_i)_{i\in I}\in (\mathbf{A}[(\mathbf{X}_j)_{j\in J}])^I\) and \(\mathbf{g} = (g_j)_{j\in J}\in (\mathbf{A}[(\mathbf{Y}_k)_{k\in K}])^J\) , we denote by \(\mathbf{f}\circ \mathbf{g}\) or \(f(\mathbf{g})\) , the family of polynomials \((f_i(\mathbf{g}))_{i\in I}\in (\mathbf{A}[(\mathbf{Y}_k)_{k\in K}])^I\) . If we denote by \(\tilde{\mathbf{f}}\) the mapping \(\mathbf{x} \mapsto (f_i(\mathbf{x}))_{i \in I}\) of \(\mathrm{E}^{\prime J}\) into \(\mathrm{E}^{\prime I}\) (where \(\mathrm{E}'\) is a unital associative and commutative A-algebra), then the relation (2) implies

\[
(\mathbf {f} \circ \mathbf {g}) ^ {\sim} = \tilde {\mathbf {f}} \circ \tilde {\mathbf {g}}.\tag{3}
\]

If \(\mathbf{h} = (h_k)_{k\in K}\in (\mathbf{A}[(\mathbf{Z}_l)_{l\in L}])^{\mathbf{K}}\) , it follows from (2) that :

\[
\mathbf {f} \circ (\mathbf {g} \circ \mathbf {h}) = (\mathbf {f} \circ \mathbf {g}) \circ \mathbf {h}.\tag{4}
\]

PROPOSITION 2. --- Let \(\mathbf{a} = (a_i)_{i \in I}\) be a family of elements of A and let \(u \in A[X]\) . If v is the polynomial obtained by substituting \(X_i + a_i\) for \(X_i\) for each \(i \in I\) , then the constant term of v is equal to \(u(\mathbf{a})\) .

The constant term of v is obtained by substituting 0 for \(X_{i}\) in v for each \(i \in I\) . The result therefore follows by (2).

COROLLARY 1. --- Let m be the ideal of polynomials \(u \in A[X]\) such that \(u(\mathbf{a}) = 0\) . Then m is generated by the polynomials \(X_{i} - a_{i}\) (for \(i \in I\) ).

It is clear that \(X_{i} - a_{i} \in m\) for each \(i \in I\) . Let \(u \in m\) and let v be as in Prop. 2. Since v has no constant term, there exists a family \((\mathbf{P}_{i})_{i \in I}\) of polynomials in A{[}X{]} with finite support such that

\[
v (\mathbf {X}) = \sum_ {i \in I} \mathrm{X} _ {i} \cdot \mathrm{P} _ {i} (\mathbf {X}).
\]

If we replace \(X_{i}\) by \(X_{i}-a_{i}\) for each \(i\in I\) in the above equation, we obtain a relation of the form \(u(\mathbf{X})=\sum_{i\in I}\left(\mathbf{X}_{i}-a_{i}\right)\cdot\mathbf{P}_{i}^{\prime}(\mathbf{X})\) , whence the corollary.

COROLLARY 2. --- Let \(\mathbf{X} = (\mathbf{X}_i)_{i \in I}\) and \(\mathbf{Y} = (\mathbf{Y}_i)_{i \in I}\) be two families of indeterminates. The set of polynomials \(u \in \mathbf{A}[\mathbf{X}, \mathbf{Y}]\) such that \(u(\mathbf{X}, \mathbf{X}) = 0\) is the ideal of \(\mathbf{A}[\mathbf{X}, \mathbf{Y}]\) generated by the polynomials \(\mathbf{X}_i - \mathbf{Y}_i\) (for \(i \in I\) ).

This corollary results directly from Cor. 1 on replacing A by A{[}Y{]} and \(a_{i}\) by \(Y_{i}\) , interpreting A{[}X, Y{]} as polynomial ring in the \(X_{i}\) with coefficients in A{[}Y{]}.

PROPOSITION 3. --- Let \(u \in \mathbf{A}[\mathbf{X}]\) and let \(\mathbf{X} \cdot \mathbf{Z}\) be the family \((\mathbf{X}_i\mathbf{Z})_{i \in I}\) of elements of the polynomial ring \(\mathbf{A}[\mathbf{X}][\mathbf{Z}]\) . The coefficient of \(\mathbf{Z}^k\) in \(u(\mathbf{X}, \mathbf{Z})\) is the homogeneous component of degree \(k\) of \(u\) , for every positive integer \(k\) .

It suffices to prove this Prop. in the case where \(u\) is a monomial, and in this case the result is clear.

COROLLARY. --- For a polynomial \(u \in A[X]\) to be homogeneous of degree \(k\) it is necessary and sufficient that :

\[
u (\mathbf {X} \cdot \mathbf {Z}) = u (\mathbf {X}) \cdot \mathbf {Z} ^ {k}.
\]

Remark. --- Let \(x \in A^{I}\) and let f be the mapping \(u \mapsto u(x)\) of A{[}X{]} into A. Given an A-module M, we consider the homomorphism \(1 \otimes f\) of M{[}X{]} = M \(\otimes_{A} A[X]\) into \(\mathbf{M} \otimes_{\mathbf{A}} \mathbf{A} = \mathbf{M}\) . For each \(v \in \mathbf{M}[\mathbf{X}]\) we have \((1 \otimes f)(v) = v(\mathbf{x})\) . If \(v = \sum_{\nu \in \mathbf{N}^{(I)}} e_{\nu} X^{\nu}\) , then \(v(\mathbf{x}) = \sum_{\nu \in \mathbf{N}^{(I)}} \mathbf{x}^{\nu} e_{\nu}\) .

